% =============================================================================
% Sección: Caso de Negocio (nombre literal del enunciado)
% =============================================================================
\section{Caso de Negocio}

\subsection{Enunciado del caso}

El proyecto semestral plantea el desafío de construir una Arquitectura de
Solución BI para un problema real. El caso propuesto es el siguiente:

\begin{quote}
  \itshape
  «Una compañía internacional de transportes decide fusionar sus líneas de
  negocio en términos de información, para tener una vista única de sus
  clientes».
\end{quote}

La empresa es \textbf{Transportes Rapiditos S.A.}, que opera seis líneas de
negocio. Cada línea vende servicios distintos (pasajes, viajes, envíos y
pedidos), por lo que la información de un mismo cliente queda repartida entre
las líneas. El caso pide integrar esa información para analizar al cliente de
manera única, sin importar la línea en que compre.

\subsection{Líneas de negocio}

El enunciado define las siguientes líneas de negocio. La tabla indica además
el medio de transporte y el tipo de servicio principal con que cada línea se
representa en el modelo dimensional (atributos de la dimensión
\bd{dim_servicio}).

\begin{table}[H]
  \centering
  \caption{Líneas de negocio de Transportes Rapiditos S.A.}
  \label{tab:lineas_negocio}
  \begin{tabular}{@{}lll@{}}
    \toprule
    \textbf{Línea de negocio} & \textbf{Medio de transporte} & \textbf{Tipo de servicio} \\
    \midrule
    Buses InterUrbanos Rapiditos       & Bus             & Pasajeros \\
    Trenes Rapiditos                   & Tren            & Pasajeros \\
    Radio Taxis Rapiditos              & Taxi            & Pasajeros \\
    Logística y Distribución Rapiditos & Camión          & Carga \\
    Aerolíneas Rapiditos               & Avión           & Pasajeros \\
    URapid Delivery                    & Moto/bicicleta  & Carga \\
    \bottomrule
  \end{tabular}
\end{table}

\subsection{Objetivo}

El objetivo definido en el enunciado es:

\begin{quote}
  \itshape
  «Determinar qué tipos de clientes tiene la compañía y su comportamiento
  anual, en base a los últimos 3 años».
\end{quote}

El comportamiento se analiza según la compra, los medios de transporte, la
zona y la edad, entre otras variables. A partir de este objetivo, la solución
BI debe responder las siguientes preguntas de negocio:

\begin{enumerate}[leftmargin=*]
  \item ¿Qué tipos de clientes tiene la compañía (personas o empresas, por
        género, rango de edad y lugar de residencia)?
  \item ¿Cuánto compra cada tipo de cliente por año y cómo cambia ese monto en
        los últimos tres años?
  \item ¿Qué líneas de negocio y medios de transporte usa cada tipo de
        cliente, y cuántos clientes compran en más de una línea?
  \item ¿Desde qué zonas y hacia qué zonas se mueven los clientes?
  \item ¿Por qué canales compran los clientes y con qué medios de pago?
\end{enumerate}

\subsection{Consideraciones del enunciado}

\begin{itemize}[leftmargin=*]
  \item El modelamiento dimensional debe considerar un modelo tipo estrella.
  \item La reportería debe presentarse para un perfil de usuario:
        operacional/táctico o ejecutivo.
  \item Las recomendaciones deben estar asociadas al negocio del caso.
  \item Se pueden usar datasets públicos disponibles en la web.
\end{itemize}

\subsection{Alcance del Datawarehouse}

El Datawarehouse \bd{dwh_rapiditos} tiene el siguiente alcance:

\begin{itemize}[leftmargin=*]
  \item \textbf{Proceso de negocio:} compras de servicios que realizan los
        clientes en las seis líneas de negocio.
  \item \textbf{Grano:} una fila por cada transacción de compra de un cliente
        por un servicio de una línea de negocio, en una fecha. Por ejemplo, un
        pasaje de bus, tren o avión; un viaje de taxi; un envío de logística o
        un pedido de delivery.
  \item \textbf{Vista única del cliente:} cada cliente se identifica con una
        llave natural común a todas las líneas (RUT o documento), lo que
        permite consolidar sus compras en una sola dimensión
        (\bd{dim_cliente}).
  \item \textbf{Modelo:} un esquema estrella con una tabla de hechos
        (\bd{fact_venta}) y seis dimensiones: tiempo, cliente, servicio, zona
        (como origen y destino), canal de venta y medio de pago.
  \item \textbf{Período:} los últimos tres años. Los años exactos se fijan al
        definir la fuente de datos, en el Avance N°2.
  \item \textbf{Moneda:} los montos se expresan en pesos chilenos (CLP).
  \item \textbf{Motor:} MariaDB, administrado con XAMPP y phpMyAdmin.
\end{itemize}

Este Avance N°1 incluye la definición del modelo multidimensional y el
script SQL que crea su estructura, \emph{sin datos}. El poblamiento mediante
ETL, el cubo, los KPI y el dashboard de visualización se desarrollan en los
avances siguientes, de acuerdo con la planificación de la Tabla~\ref{tab:entregas}.

\begin{table}[H]
  \centering
  \caption{Planificación de entregas del proyecto}
  \label{tab:entregas}
  \begin{tabularx}{\linewidth}{@{}llX@{}}
    \toprule
    \textbf{Entrega} & \textbf{Fecha} & \textbf{Contenido} \\
    \midrule
    Avance N°1 & 05-11-2026 & Caso de negocio, modelo multidimensional
                              estrella y script SQL de creación (sin datos). \\
    Avance N°2 & 19-11-2026 & Fuentes de información, ETL y poblamiento del
                              Datawarehouse. \\
    Avance N°3 & 10-12-2026 & Cubo, KPI, dashboard de visualización,
                              recomendaciones y presentación. \\
    \bottomrule
  \end{tabularx}
\end{table}
